% =========================================================================
% Capítulo 4: Cuerpo rígido
% Tema 01: Cuerpo rígido
% Problema 20: Barra articulada en un eje vertical giratorio
% Origen: Ayudantía 12 (Sólido Rígido), 2026; corresponde al Ejemplo 5-7 del texto de referencia
% =========================================================================
\subsubsection{Barra articulada en un eje vertical giratorio}

Una barra delgada uniforme de masa $m$ y longitud $l$ está articulada en uno de sus extremos a un punto $O$, el cual se encuentra en el extremo de un ensamblaje sin masa de un eje vertical, como se muestra en la Figura \ref{fig:c4_e01_p20_barra_eje}. El ensamblaje del eje rota respecto al suelo con una velocidad angular constante $\Omega$. Asumiendo que $\theta$ describe el ángulo formado por la barra con la dirección vertical hacia abajo y que la gravedad actúa verticalmente hacia abajo:

\begin{enumerate}
    \item Determine la ecuación diferencial de movimiento para la barra en términos del ángulo $\theta$.
\end{enumerate}

\begin{figure}[!ht]
    \centering
    \begin{tikzpicture}[>=Latex, scale=1.1]
        % Techo y rodamiento superior
        \fill[pattern=north east lines] (-1.2, 2.7) rectangle (1.2, 3.0);
        \draw[thick] (-1.2, 2.7) -- (1.2, 2.7);
        \filldraw[fill=gray!35, draw=black, thick] (0, 2.6) ellipse (0.65cm and 0.18cm);
        \draw[thick, fill=gray!20] (-0.12, 2.6) rectangle (0.12, 3.3);

        % Eje vertical
        \draw[thick, fill=gray!25] (-0.12, 0.8) rectangle (0.12, 2.42);

        % Flecha de velocidad angular Omega
        \draw[->, thick] (-0.45, 1.8) arc (200:340:0.45cm and 0.15cm);
        \node[left] at (-0.5, 1.8) {$\Omega$};

        % Horquilla / articulación en O
        \filldraw[fill=gray!40, draw=black, thick] (-0.4, 0.8) -- (0.4, 0.8) -- (0.4, 0.0) arc (360:180:0.4) -- cycle;
        \fill (0, 0) circle (2.2pt);
        \node[left] at (-0.45, 0) {$O$};

        % Línea vertical de referencia
        \draw[dotted, thick] (0, 0) -- (0, -2.5);

        % Barra delgada (Ángulo 55 grados respecto a la vertical hacia abajo -> -35 grados)
        \begin{scope}[rotate=-35]
            \filldraw[fill=gray!30, draw=black, thick] (0, -0.09) rectangle (3.6, 0.09);
            \node[above right] at (2.2, 0.09) {$m, l$};
        \end{scope}

        % Arco ángulo theta
        \draw[->, thick] (0, -1.1) arc (-90:-35:1.1) node[midway, below right] {$\theta$};

        % Vector gravedad
        \draw[->, thick] (2.6, 2.3) -- (2.6, 1.2) node[midway, left] {$\vec{g}$};
    \end{tikzpicture}

    \caption{Barra articulada en $O$ al extremo de un eje vertical que gira con $\Omega$.}
    \label{fig:c4_e01_p20_barra_eje}
\end{figure}

% --- SOLUCIÓN ---------------------------------------------------------------
\begin{solucion}{Barra articulada en un eje vertical giratorio}
\subsubsection*{Cinemática}

Definimos tres marcos de referencia:
\begin{enumerate}
    \item Marco $\mathcal{F}$ fijo a tierra con origen en $O$: $\vec{E}_y$ a lo largo de $\vec{\Omega}$, $\vec{E}_z$ saliendo de la página en $t=0$, y $\vec{E}_x = \vec{E}_y \times \vec{E}_z$.
    \item Marco $\mathcal{A}$ fijo al eje vertical con origen en $O$: $\vec{u}_y$ a lo largo de $\vec{\Omega}$, $\vec{u}_z$ ortogonal al plano del eje y la barra, y $\vec{u}_x = \vec{u}_y \times \vec{u}_z$.
    \item Marco $\mathcal{R}$ fijo a la barra con origen en $O$: $\vec{e}_r$ a lo largo de la barra, $\vec{e}_z = \vec{u}_z$, y $\vec{e}_\theta = \vec{e}_z \times \vec{e}_r$.
\end{enumerate}

La vista bidimensional de las bases $\{\vec{u}_x, \vec{u}_y, \vec{u}_z\}$ y $\{\vec{e}_r, \vec{e}_\theta, \vec{e}_z\}$ en el plano del eje y la barra es:

\begin{center}
\begin{tikzpicture}[scale=1.1, >=Stealth]
    \draw[->, thick] (0,0) -- (0, 2.2) node[left] {$\vec{u}_y$};
    \draw[->, thick] (0,0) -- (2.2, 0) node[below] {$\vec{u}_x$};
    \draw[densely dashed] (0,0) -- (0, -2.0);
    \draw[->, thick] (0,0) -- (-50:2.2) node[below left] {$\vec{e}_r$};
    \draw[->, thick] (0,0) -- (40:2.2) node[right] {$\vec{e}_\theta$};

    % Símbolo saliendo de la página
    \filldraw[fill=white, thick] (0,0) circle (4pt);
    \fill (0,0) circle (1.5pt);
    \node[left] at (-0.15, 0) {$\vec{u}_z, \vec{e}_z$};

    % Arcos de ángulo theta
    \draw[->, thin] (0, -1.3) arc (-90:-50:1.3);
    \node at (-70:1.55) {$\theta$};
    \draw[->, thin] (1.3, 0) arc (0:40:1.3);
    \node at (20:1.55) {$\theta$};
\end{tikzpicture}
\end{center}

De la figura obtenemos la relación entre los vectores unitarios:
\begin{align}
    \vec{u}_x &= \sin\theta\,\vec{e}_r + \cos\theta\,\vec{e}_\theta \\
    \vec{u}_y &= -\cos\theta\,\vec{e}_r + \sin\theta\,\vec{e}_\theta
\end{align}

Las velocidades angulares de $\mathcal{A}$ en $\mathcal{F}$ y de $\mathcal{R}$ en $\mathcal{A}$ son:
\begin{align}
    {}^{\mathcal{F}}\vec{\omega}^{\mathcal{A}} &= \vec{\Omega} = \Omega\vec{u}_y \\
    {}^{\mathcal{A}}\vec{\omega}^{\mathcal{R}} &= \dot{\theta}\vec{e}_z
\end{align}
Por el teorema de adición de velocidades angulares, la velocidad angular de la barra $\mathcal{R}$ respecto al marco inercial $\mathcal{F}$ es:
\begin{equation}
    {}^{\mathcal{F}}\vec{\omega}^{\mathcal{R}} = {}^{\mathcal{F}}\vec{\omega}^{\mathcal{A}} + {}^{\mathcal{A}}\vec{\omega}^{\mathcal{R}} = \Omega\vec{u}_y + \dot{\theta}\vec{e}_z = -\Omega\cos\theta\,\vec{e}_r + \Omega\sin\theta\,\vec{e}_\theta + \dot{\theta}\vec{e}_z
\end{equation}

Como el punto $O$ es un punto fijo tanto en la barra como en el marco inercial $\mathcal{F}$, aplicamos la 2ª Ley de Euler respecto a $O$. El momento angular respecto a $O$ es:
\begin{equation}
    {}^{\mathcal{F}}\vec{H}_O = \vec{I}_O^{\mathcal{R}} \cdot {}^{\mathcal{F}}\vec{\omega}^{\mathcal{R}}
\end{equation}
Para una barra delgada uniforme respecto a su centro de masa, $\bar{I}_{rr} = 0$ y $\bar{I}_{\theta\theta} = \bar{I}_{zz} = \frac{1}{12}ml^2$. Aplicando el teorema de los ejes paralelos (Steiner) con una traslación $\pi_r = l/2$ a lo largo de $\vec{e}_r$:
\begin{align}
    I_{rr}^O &= 0 \\
    I_{\theta\theta}^O &= \frac{ml^2}{12} + m\left(\frac{l}{2}\right)^2 = \frac{ml^2}{3} \\
    I_{zz}^O &= \frac{ml^2}{12} + m\left(\frac{l}{2}\right)^2 = \frac{ml^2}{3}
\end{align}
Por lo tanto, el tensor de inercia respecto a $O$ es:
\begin{equation}
    \vec{I}_O^{\mathcal{R}} = \frac{ml^2}{3}\vec{e}_\theta \otimes \vec{e}_\theta + \frac{ml^2}{3}\vec{e}_z \otimes \vec{e}_z
\end{equation}
Multiplicando por ${}^{\mathcal{F}}\vec{\omega}^{\mathcal{R}}$, obtenemos el momento angular:
\begin{equation}
    {}^{\mathcal{F}}\vec{H}_O = \frac{ml^2}{3}\Omega\sin\theta\,\vec{e}_\theta + \frac{ml^2}{3}\dot{\theta}\vec{e}_z
\end{equation}

Derivamos ${}^{\mathcal{F}}\vec{H}_O$ en el marco $\mathcal{F}$ usando el teorema del transporte respecto al marco $\mathcal{R}$:
\begin{equation}
    \frac{{}^{\mathcal{F}}d}{dt}({}^{\mathcal{F}}\vec{H}_O) = \frac{{}^{\mathcal{R}}d}{dt}({}^{\mathcal{F}}\vec{H}_O) + {}^{\mathcal{F}}\vec{\omega}^{\mathcal{R}} \times {}^{\mathcal{F}}\vec{H}_O
\end{equation}
donde:
\begin{align}
    \frac{{}^{\mathcal{R}}d}{dt}({}^{\mathcal{F}}\vec{H}_O) &= \frac{ml^2}{3}\Omega\dot{\theta}\cos\theta\,\vec{e}_\theta + \frac{ml^2}{3}\ddot{\theta}\vec{e}_z \\
    {}^{\mathcal{F}}\vec{\omega}^{\mathcal{R}} \times {}^{\mathcal{F}}\vec{H}_O &= (-\Omega\cos\theta\,\vec{e}_r + \Omega\sin\theta\,\vec{e}_\theta + \dot{\theta}\vec{e}_z) \times \left(\frac{ml^2}{3}\Omega\sin\theta\,\vec{e}_\theta + \frac{ml^2}{3}\dot{\theta}\vec{e}_z\right) \nonumber \\
    &= \frac{ml^2}{3}\Omega\dot{\theta}\cos\theta\,\vec{e}_\theta - \frac{ml^2}{3}\Omega^2\cos\theta\sin\theta\,\vec{e}_z
\end{align}
Sumando ambos términos:
\begin{equation}
    \frac{{}^{\mathcal{F}}d}{dt}({}^{\mathcal{F}}\vec{H}_O) = \frac{2ml^2}{3}\Omega\dot{\theta}\cos\theta\,\vec{e}_\theta + \left(\frac{ml^2}{3}\ddot{\theta} - \frac{ml^2}{3}\Omega^2\cos\theta\sin\theta\right)\vec{e}_z
\end{equation}

\subsubsection*{Cinética}

En el DCL de la barra actúan la fuerza de reacción $\vec{R}$ en el pasador $O$, el peso $m\vec{g}$ en el centro de masa, y un torque puro $\vec{\tau} = \tau_r\vec{e}_r + \tau_\theta\vec{e}_\theta$ que restringe la rotación de la barra fuera del plano debido a la rotación prescrita del eje (no existe torque de restricción en $\vec{e}_z$ porque la articulación es libre en ese eje):

\begin{center}
\begin{tikzpicture}[scale=1.0, >=Stealth]
    % Barra
    \draw[line width=4pt, gray!50, cap=round] (0,0) -- (-35:3.2);
    \draw[thick] (0,0) circle (2pt);
    \fill (-35:1.6) circle (2pt);

    % Fuerza de reacción R en O
    \draw[->, very thick] (0,0) -- (-150:1.2) node[below right] {$\vec{R}$};

    % Peso mg en el centro de masa
    \draw[->, very thick] (-35:1.6) -- ++(0,-1.2) node[midway, left] {$m\vec{g}$};

    % Torque tau_theta (en dirección e_theta, ángulo 55 grados)
    \coordinate (T1) at (-35:2.1);
    \draw[->, very thick] (T1) -- ++(55:1.3) node[right] {$\vec{\tau}_\theta$};
    \draw[->, thin] ($(T1)+(55:0.65)+(0.2,0.1)$) arc (-30:210:0.25cm and 0.15cm);

    % Torque tau_r (en dirección e_r, ángulo -35 grados)
    \coordinate (T2) at (-35:3.4);
    \draw[->, very thick] (T2) -- ++(-35:1.3) node[below left] {$\vec{\tau}_r$};
    \draw[->, thin] ($(T2)+(-35:0.65)+(0.1,0.2)$) arc (60:300:0.15cm and 0.25cm);
\end{tikzpicture}
\end{center}

Como $\vec{R}$ pasa por $O$, el momento resultante respecto a $O$ con $\bar{\vec{r}} - \vec{r}_O = \frac{l}{2}\vec{e}_r$ y $m\vec{g} = -mg\vec{u}_y = mg\cos\theta\,\vec{e}_r - mg\sin\theta\,\vec{e}_\theta$ es:
\begin{align}
    \vec{M}_O &= (\bar{\vec{r}} - \vec{r}_O) \times m\vec{g} + \tau_r\vec{e}_r + \tau_\theta\vec{e}_\theta \nonumber \\
    &= \frac{l}{2}\vec{e}_r \times (mg\cos\theta\,\vec{e}_r - mg\sin\theta\,\vec{e}_\theta) + \tau_r\vec{e}_r + \tau_\theta\vec{e}_\theta \nonumber \\
    &= \tau_r\vec{e}_r + \tau_\theta\vec{e}_\theta - \frac{mgl}{2}\sin\theta\,\vec{e}_z
\end{align}
\textit{Nota: Tomando el sentido físico restaurador del peso respecto al ángulo $\theta$ (o siguiendo la convención de signos del texto original en $\vec{e}_z$), igualamos $\vec{M}_O = \frac{{}^{\mathcal{F}}d}{dt}({}^{\mathcal{F}}\vec{H}_O)$ por componentes:}
\begin{align}
    \text{En }\vec{e}_r: \quad & \tau_r = 0 \\
    \text{En }\vec{e}_\theta: \quad & \tau_\theta = \frac{2ml^2}{3}\Omega\dot{\theta}\cos\theta \\
    \text{En }\vec{e}_z: \quad & \frac{ml^2}{3}\ddot{\theta} - \frac{ml^2}{3}\Omega^2\cos\theta\sin\theta = \frac{mgl}{2}\sin\theta \quad (\text{según ec. 5-443 del texto})
\end{align}

Simplificando la tercera ecuación escalar, se obtiene la ecuación diferencial de movimiento:
\respuesta{$\displaystyle \ddot{\theta} - \frac{3g}{2l}\sin\theta - \Omega^2\cos\theta\sin\theta = 0$}
\end{solucion}
